\documentclass[letterpaper,10pt]{article}

\usepackage{latexsym}
\usepackage[empty]{fullpage}
\usepackage{titlesec}
\usepackage{marvosym}
\usepackage[usenames,dvipsnames]{color}
\usepackage{verbatim}
\usepackage{enumitem}
%\usepackage{xcolor}
\usepackage{color}% http://ctan.org/pkg/color

\usepackage[pdftex,colorlinks=false,pdfborderstyle={/S/U/W 1},urlbordercolor={0 0 0}]{hyperref}
\usepackage{fancyhdr}
\usepackage[italian,ngerman,english]{babel}
\usepackage{tabularx}
% only for pdflatex
% \input{glyphtounicode}

% fontawesome
\usepackage{fontawesome5}

% fixed width
\usepackage[scale=0.90,lf]{FiraMono}

% light-grey
\definecolor{light-grey}{gray}{0.83}
\definecolor{dark-grey}{gray}{0.3}
\definecolor{text-grey}{gray}{.08}

\DeclareRobustCommand{\ebseries}{\fontseries{eb}\selectfont}
\DeclareTextFontCommand{\texteb}{\ebseries}

% custom underilne
\usepackage{contour}
\usepackage[normalem]{ulem}
\renewcommand{\ULdepth}{1.8pt}
\contourlength{0.8pt}
\newcommand{\myuline}[1]{%
  \uline{\phantom{#1}}%
  \llap{\contour{white}{#1}}%
}


% font: arial
\usepackage{helvet}
\renewcommand{\familydefault}{\sfdefault}
% font: helvetica-style
%\usepackage{tgheros}
%\renewcommand*\familydefault{\sfdefault} 
%% Only if the base font of the document is to be sans serif
\usepackage[T1]{fontenc}


\pagestyle{fancy}
\fancyhf{} % clear all header and footer fields
\fancyfoot{}
\renewcommand{\headrulewidth}{0pt}
\renewcommand{\footrulewidth}{0pt}

% Adjust margins
\addtolength{\oddsidemargin}{-0.5in}
\addtolength{\evensidemargin}{0in}
\addtolength{\textwidth}{1in}
\addtolength{\topmargin}{-.8in}
\addtolength{\textheight}{1.5in}

\urlstyle{same}

\raggedbottom
\raggedright
\setlength{\tabcolsep}{0in}

% Sections formatting - serif
% \titleformat{\section}{
%   \vspace{2pt} \scshape \raggedright\large % header section
% }{}{0em}{}[\color{black} \titlerule \vspace{-5pt}]

% TODO EBSERIES
% sans serif sections
\titleformat {\section}{
    \bfseries \vspace{2pt} \raggedright \large % header section
}{}{0em}{}[\color{light-grey} {\titlerule[2pt]} \vspace{-4pt}]

% only for pdflatex
% Ensure that generate pdf is machine readable/ATS parsable
% \pdfgentounicode=1

%-------------------------
% Custom commands
\newcommand{\resumeItem}[1]{
  \item\small{
    {#1 \vspace{-1pt}}
  }
}

\newcommand{\resumeSubheading}[4]{
  \vspace{-1pt}\item
    \begin{tabular*}{\textwidth}[t]{l@{\extracolsep{\fill}}r}
      \textbf{#1} & {\color{dark-grey}\small #2}\vspace{1pt}\\ % top row of resume entry
      \textit{#3} & {\color{dark-grey} \small #4}\\ % second row of resume entry
    \end{tabular*}\vspace{-4pt}
}

\newcommand{\resumeSubSubheading}[2]{
    \item
    \begin{tabular*}{\textwidth}{l@{\extracolsep{\fill}}r}
      \textit{\small#1} & \textit{\small #2} \\
    \end{tabular*}\vspace{-7pt}
}

\newcommand{\resumeProjectHeading}[2]{
    \item
    \begin{tabular*}{\textwidth}{l@{\extracolsep{\fill}}r}
      #1 & {\color{dark-grey}} \\
    \end{tabular*}\vspace{-4pt}
}

\newcommand{\resumeSubItem}[1]{\resumeItem{#1}\vspace{-4pt}}

\renewcommand\labelitemii{$\vcenter{\hbox{\tiny$\bullet$}}$}

% CHANGED default leftmargin  0.15 in
\newcommand{\resumeSubHeadingListStart}{\begin{itemize}[leftmargin=0in, label={}]}
\newcommand{\resumeSubHeadingListEnd}{\end{itemize}}
\newcommand{\resumeItemListStart}{\begin{itemize}}
\newcommand{\resumeItemListEnd}{\end{itemize}}

\color{text-grey}

%-------------------------------------------
%%%%%%  RESUME STARTS HERE  %%%%%%%%%%%%%%%%%%%%%%%%%%%%


\begin{document}
\selectlanguage{german}
%----------HEADING----------
% See icons: http://mirrors.ibiblio.org/CTAN/fonts/fontawesome/doc/fontawesome.pdf
\begin{center}
    \textbf{\Huge Niccol\`o Zuppichini} \\ \vspace{5pt}
    \small \faPhone* \texttt{+41 76 275 07 06} \hspace{1pt} $|$
    \hspace{1pt} \faEnvelope \hspace{2pt} \texttt{niccolo.zuppichini@gmail.com} \hspace{1pt}  \\
    \hspace{1pt} \faLinkedin \hspace{2pt} \texttt{\href{https://www.linkedin.com/in/niccolo-zuppichini}{www.linkedin.com/in/niccolo-zuppichini}} \hspace{1pt} \\
%    \hspace{1pt} \faGithub \hspace{2pt} \texttt{rage997} \hspace{1pt} $|$ \\
    \hspace{1pt} \faHome \hspace{2pt}\texttt{\href{niccolozuppichini.com}{niccolozuppichini.com}}   \hspace{1pt} $|$
    \hspace{1pt} \faMapMarker* \hspace{2pt}\texttt{Bellinzona, Schweiz}
    \\ \vspace{-3pt}
\end{center}

%-----------EXPERIENCE-----------
\section{BERUFSERFAHRUNG}
  \resumeSubHeadingListStart
  
%-----------Luminocity-----------
\resumeSubheading
    {\href{https://luminocity.ch/}{Luminocity AG}}{April 2025 -- April 2026}
    {Senior Software Engineer}{Schaffhausen, CH (Remote)}
    \resumeItemListStart
      \resumeItem{Beitrag zur Kernerwicklung der interaktiven 3D-Plattform Lemonate, mit Echtzeit-Rendering und Deployment komplexer browserbasierter Anwendungen.}
      \resumeItem{Entwicklung zentraler Rendering-Engine-Funktionen, darunter Echtzeit-Partikelsysteme, Post-Processing-Pipelines und Low-Level-Grafikarchitektur.}
      \resumeItem{Implementierung fortgeschrittener Rendering-Techniken und Performance-Optimierungen in einer WebGL-basierten Engine f\"ur plattform\"ubergreifende 3D-Visualisierung.}
      \resumeItem{F\"uhrung der Engine-Architektur-Entscheidungen und Entwicklung von Tools zur Optimierung der 3D-Workflows in Web-Umgebungen.}
    \resumeItemListEnd

%-----------Mindcraft AG-----------
\resumeSubheading
    {\href{https://mndcr.com/}{Mindcraft AG}}{Oktober 2022 -- April 2025}
    {Senior Software Engineer}{Bern, CH (Remote)}
    \resumeItemListStart
      \resumeItem{Konzeption und Implementierung einer leistungsf\"ahigen Web-Rendering-Engine f\"ur die Meshvalley-Plattform.}
      \resumeItem{Aufbau von End-to-End-3D-Scanning-Pipelines, einschlie\ss{lich} Asset-Verarbeitung, -Optimierung und -Vorbereitung f\"ur Echtzeit-Rendering.}
      \resumeItem{Architektur von Systemen f\"ur Hosting, Streaming und interaktive Visualisierung gro\ss{skaliger} 3D-Datens\"atze im Browser.}
      \resumeItem{Entwicklung geometrischer Algorithmen und Datenverarbeitungs-Workflows zur effizienten Bearbeitung komplexer gescannter Umgebungen.}
    \resumeItemListEnd

%-----------Panua-----------
    \resumeSubheading
      {\href{https://panua.ch/}{Panua}}{Juli 2022 -- November 2023}
      {Web Developer / DevOps Engineer (Freelance)}{Lugano, CH (Remote)}
      \resumeItemListStart
      \resumeItem{Zusammenarbeit mit dem Panua-Team auf Stundenbasis, mit Expertise in Web-Entwicklung, DevOps und Software Engineering.}
      \resumeItem{Implementierung einer Software zur Generierung, Validierung und Authentifizierung von Maschinenfingerabdr\"ucken der Nutzer, um Software-Piraterie bei den Produkten von Panua Technologies zu verhindern und die Produktsicherheit zu erh\"ohen.}
      \resumeItem{Gew\"ahrleistung hoher Verf\"ugbarkeit und Leistung der Softwareplattformen von Panua.}
      \resumeItem{Entwicklung automatisierter Deployment- und Test-Pipelines mit Docker, zur Verbesserung der Effizienz und Zuverl\"assigkeit von Software-Releases.}
    \resumeItemListEnd

%-----------HEGIAS-----------
    \resumeSubheading
      {\href{https://hegias.com/en/}{Hegias}}{September 2018 -- September 2020}
      {Software Engineer | DevOps Engineer}{Lugano \& Z\"urich, CH (Hybrid)}
      \resumeItemListStart
       \resumeItem{F\"uhrung eines Agile-Teams aus 5 Entwicklern hin zu Continuous Deployment und 20\% h\"oherer Produktivit\"at durch Pair Programming.}
        \resumeItem{Verantwortlich f\"ur die Unix-Server-Infrastruktur, die Anfragen von Tausenden von Kunden verarbeitet.}
        \resumeItem{Selbstst\"andige L\"osung der Performance-Probleme der Rendering-Engine durch fortgeschrittene Rendering-Techniken, 3D-Objektkompression und maschinelles Lernen zur 3D-Objektklassifizierung.}
        \resumeItem{Reduzierung der Laufzeit der 3D-Objekt-Verarbeitungspipeline um \"uber 250\% durch Verteilung der Server-Infrastruktur auf verschiedene Knoten.}
        \resumeItem{Das Unternehmen erhielt \"uber 2,1 Millionen Schweizer Franken an Finanzierung und wurde zu einer der \href{https://www.eu-startups.com/2019/02/10-swiss-startups-to-watch-in-2019/}{10 vielversprechendsten Startups} ausgew\"ahlt.}
      \resumeItemListEnd
      
%-----------CS-----------
    \resumeSubheading
      {\href{https://www.ics.usi.ch/}{Institut f\"ur Computational Science}}{Juni 2018 -- September 2018}
      {Junior Researcher}{Lugano, CH}
      \resumeItemListStart
        \resumeItem{Beitrag zur Entwicklung von \href{https://bitbucket.org/zulianp/utopia/src/master/}{Utopia}, einer leistungsf\"ahigen C++-Parallel-Bibliothek zur L\"osung nichtlinearer, hochdimensionaler Probleme.}
        \resumeItem{Vollst\"andiger Quellcode verf\"ugbar im \href{https://bitbucket.org/zulianp/utopia/src/master/}{Repository des Projekt-Warters}}
    \resumeItemListEnd

  \resumeSubHeadingListEnd


% --------- Education -----------
\section {AUSBILDUNG}
\resumeSubHeadingListStart
   \resumeSubheading
     {Universit\`a della Svizzera Italiana}{September 2019 -- Juni 2022}
     {Master of Science in Computational Science}{Lugano, CH}
       \resumeItemListStart
       \resumeItem 
           {\textbf{Abschlussarbeit}: \href{https://github.com/Rage997/MasterThesis}{Modellierung invasiver Arten mittels relationaler Ereignismodelle.}}
     \resumeItem {\textbf{F\"acher}: Stochastische Optimierung, Numerische Analysis, PDE/ODE, Deep Learning.}
     \resumeItem{\textbf{Notendurchschnitt}: 7.66}
       \resumeItemListEnd
 \resumeSubHeadingListEnd
\resumeSubHeadingListStart
   \resumeSubheading
     {Universit\`a della Svizzera Italiana}{September 2016 -- Juni 2019}
     {Bachelor of Science in Informatik}{Lugano, CH}
       \resumeItemListStart
       \resumeItem 
           {\textbf{Projekt}: \href{https://github.com/Rage997/BachelorProject}{Eine robuste Pipeline zur Verarbeitung von 3D-Objekten.}}
     \resumeItem {\textbf{F\"acher}: Datenstrukturen, Algorithmen, Datenbanken, Computersysteme, Machine Learning.}
      \resumeItem{\textbf{Notendurchschnitt}: 7.45}
       \resumeItemListEnd
 \resumeSubHeadingListEnd


%-----------PROJECTS-----------

\section{PROJEKTE}
    \resumeSubHeadingListStart
     \resumeProjectHeading
         {\textbf{3D-DICOM-Viewer}} {2022}
         \resumeItemListStart
           \resumeItem{Ein moderner, browserbasierter Medizinbild-Viewer zur Darstellung von DICOM-Dateien in 2D und 3D. Kostenlos und Open Source.}
           \resumeItem{Mit Fokus auf Datenschutz entwickelt; die gesamte Verarbeitung erfolgt zu 100\% clientseitig, ohne Server-Uploads.}
           \resumeItem{Als Lernwerkzeug entwickelt, um DICOM-Parsing, volumetrisches Rendering und GPU-beschleunigte Visualisierungstechniken im Browser zu erkunden.}
         \resumeItemListEnd
          
        \resumeProjectHeading
          {\textbf{3D-Spektralsegmentierung}}{2021}
          \resumeItemListStart
            \resumeItem{Konzeption und Implementierung eines \href{https://www.blender.org}{Blender}-Add-ons zur automatischen Erkennung von 3D-Objektmerekmalen mit einem eigenen Spektralsegmentierungs-Algorithmus.}
            \resumeItem{Das Add-on wurde \"uber 200-mal heruntergeladen und erhielt viel positives Feedback.}
			\resumeItem{Der vollst\"andige Quellcode ist auf meinem \href{https://github.com/Rage997/Spectral-Clustering-mesh-segmentation}{GitHub-Profil} zu finden.}

          \resumeItemListEnd
    \resumeSubHeadingListEnd


%
%-----------PROGRAMMING SKILLS-----------
\section{KOMPETENZEN}
\newcolumntype{L}{>{\raggedleft}p{0.14\textwidth}}
\newcolumntype{R}{p{0.8\textwidth}}

\begin{minipage}[ht]{0.48\textwidth}
\textbf{Programmiersprachen:} \\
Python C++, C\#, OO \\ Javascript, Rust
\vspace{1em}

\textbf{Frameworks:} \\
Numpy, Pandas, Scikit-learn, PyTorch, Cuda, OpenMPI, \\ 
NodeJS, mongoDB, WebGL, ThreeJS
\end{minipage}
\begin{minipage}[ht]{0.48\textwidth}
\textbf{Deployment:} \\
AWS, Kubernetes, Rancher, Docker, Linux

\vspace{1em}
\textbf{Sprachen:} \\
Italienisch: Muttersprache \\
Englisch: Vollst\"andige Berufskompetenz\\
Spanisch: Konversationssichere Kenntnisse \\
Deutsch: Grundkenntnisse
%Franz\"osisch: Berufliche Kompetenz

\end{minipage}


%-------------------------------------------
\end{document}
